\documentclass[UTF8,10pt,letterpaper,fontset=fandol]{ctexart}
\usepackage[margin=0.9in]{geometry}
\usepackage{amsmath,amssymb,enumitem,needspace}
\usepackage[unicode,hidelinks]{hyperref}
\pagestyle{plain}
\raggedbottom
\setlength{\parindent}{0pt}
\setlength{\parskip}{1pt}
\setlength{\emergencystretch}{2em}
\linespread{1.06}
\AtBeginDocument{\setlength{\abovedisplayskip}{5pt}\setlength{\belowdisplayskip}{5pt}\setlength{\abovedisplayshortskip}{3pt}\setlength{\belowdisplayshortskip}{3pt}}
\setcounter{secnumdepth}{0}
\setlist[enumerate]{label=(\arabic*),leftmargin=2em,itemsep=0.2em,topsep=0.2em}
\ctexset{section={format=\Large\bfseries,beforeskip=1.1ex,afterskip=0.7ex}}
\newcommand{\C}{\mathbb C}
\newcommand{\R}{\mathbb R}
\newcommand{\F}{\mathbb F}
\newcommand{\ip}[2]{\langle #1,#2\rangle}
\DeclareMathOperator{\rank}{rank}
\DeclareMathOperator{\tr}{tr}
\DeclareMathOperator{\im}{im}
\begin{document}
\begin{center}
{\LARGE 作业 9}\\[0.5em]
2026 年 11 月 4 日\\[0.25em]
姓名：YuanHao Tang\quad 学号：3120265621
\end{center}
\section{第 1 题}
设 $\mathcal A\in L(\C_n[\lambda])$，$\mathcal A f(\lambda)=f(\lambda+a)$，$a\ne0$。利用作业 5 第 3 题中基 $1,\lambda,\ldots,\lambda^n$ 下的矩阵 $A$，求 $A$ 的最小多项式。
\par\Needspace{4\baselineskip}\textbf{解答}\par
\[
 A_{ij}=\begin{cases}\binom ji a^{j-i},&0\le i\le j\le n,\\0,&i>j,\end{cases}
 \qquad \Delta=\mathcal A-I.
\]
\[
 \Delta(\lambda^k)=(\lambda+a)^k-\lambda^k
 =ka\lambda^{k-1}+\text{低次项},\qquad k\ge1.
\]
\[
 \Delta^{n+1}=0,\qquad \Delta^n(\lambda^n)=n!a^n\ne0
 \quad\Longrightarrow\quad (A-I)^{n+1}=0,\quad(A-I)^n\ne0.
\]
\[
 \boxed{m_A(z)=(z-1)^{n+1}}.
\]
\section{第 2 题}
在内积空间 $V$ 中令 $\|\alpha\|=\sqrt{\ip\alpha\alpha}$。证明：(1) 三角不等式；(2) Cauchy--Schwarz 不等式（定理 4.3）。
\par\Needspace{4\baselineskip}\textbf{解答}\par
沿用课件约定：内积对第一变量共轭线性、对第二变量线性。先证 (2)。
\[
 \beta=0\ \Longrightarrow\ |\ip\alpha\beta|=0.
\]
当 $\beta\ne0$ 时，取 $c=\ip\beta\alpha/\|\beta\|^2$，则
\begin{align*}
 0&\le\|\alpha-c\beta\|^2\\
  &=\|\alpha\|^2-c\ip\alpha\beta-\bar c\ip\beta\alpha+|c|^2\|\beta\|^2\\
  &=\|\alpha\|^2-\frac{|\ip\alpha\beta|^2}{\|\beta\|^2}.
\end{align*}
\[
 \boxed{|\ip\alpha\beta|\le\|\alpha\|\,\|\beta\|}.
\]
\textbf{(1)}
\begin{align*}
 \|\alpha+\beta\|^2
 &=\|\alpha\|^2+2\operatorname{Re}\ip\alpha\beta+\|\beta\|^2\\
 &\le\|\alpha\|^2+2|\ip\alpha\beta|+\|\beta\|^2
 \le(\|\alpha\|+\|\beta\|)^2.
\end{align*}
\[
 \boxed{\|\alpha+\beta\|\le\|\alpha\|+\|\beta\|}.
\]
\newpage
\section{第 3 题}
令 $V=\R_n[\lambda]$，$\ip pq=\int_{-1}^1p(\lambda)q(\lambda)\,d\lambda$。证明其为内积，求基 $1,\lambda,\ldots,\lambda^n$ 的 Gram 矩阵，并用 Gram--Schmidt 过程求正交规范基 $\alpha_kL_k(\lambda)$，$0\le k\le n$。
\par\Needspace{4\baselineskip}\textbf{解答}\par
\textbf{(1)}
\[
 \ip pq=\ip qp,\qquad
 \ip{ap+bq}r=a\ip pr+b\ip qr\quad(a,b\in\R),
\]
\[
 \ip pp=\int_{-1}^1p^2\,d\lambda\ge0,\qquad
 \int_{-1}^1p^2\,d\lambda=0\ \Longrightarrow\ p|_{[-1,1]}=0\ \Longrightarrow\ p=0.
\]
\textbf{(2)}
\[
 \boxed{G=(g_{ij})_{i,j=0}^n,\qquad
 g_{ij}=\int_{-1}^1\lambda^{i+j}\,d\lambda
 =\begin{cases}\dfrac{2}{i+j+1},&i+j\text{ 为偶数},\\[3pt]0,&i+j\text{ 为奇数} .\end{cases}}
\]
\textbf{(3)}
\[
 q_0=1,\qquad
 q_k=\lambda^k-\sum_{j=0}^{k-1}\frac{\ip{q_j}{\lambda^k}}{\ip{q_j}{q_j}}q_j,\qquad
 e_k=\frac{q_k}{\|q_k\|}.
\]
\begin{align*}
 q_1&=\lambda,\qquad
 q_2=\lambda^2-\frac{2/3}{2}=\lambda^2-\frac13,\\
 q_3&=\lambda^3-\frac{2/5}{2/3}\lambda=\lambda^3-\frac35\lambda.
\end{align*}
取标准 Legendre 多项式（归一化为 $L_k(1)=1$）：
\[
 L_k(\lambda)=\frac{1}{2^kk!}\frac{d^k}{d\lambda^k}(\lambda^2-1)^k,
 \qquad c_k=\frac{(2k)!}{2^k(k!)^2}.
\]
分部积分 $k$ 次；端点处的边界项均为零。
\[
 \ip{L_k}{\lambda^j}
 =\frac{(-1)^k}{2^kk!}\int_{-1}^1(\lambda^2-1)^k
 \frac{d^k\lambda^j}{d\lambda^k}\,d\lambda=0\quad(0\le j<k).
\]
\[
 L_k=c_k\lambda^k+\cdots
 \quad\Longrightarrow\quad q_k=\frac{L_k}{c_k}.
\]
令 $I_k=\int_{-1}^1(1-\lambda^2)^k\,d\lambda$。
\[
 I_0=2,\qquad
 0=\int_{-1}^1\frac{d}{d\lambda}\bigl[\lambda(1-\lambda^2)^k\bigr]d\lambda
 =(2k+1)I_k-2kI_{k-1}\quad(k\ge1),
\]
\[
 I_k=\frac{2^{2k+1}(k!)^2}{(2k+1)!}.\]
\[
 \|L_k\|^2=\frac{(-1)^k}{2^kk!}\int_{-1}^1L_k^{(k)}(\lambda)(\lambda^2-1)^k\,d\lambda=\frac{c_k}{2^k}I_k=\frac{2}{2k+1}.
\]
\[
 \boxed{\alpha_k=\sqrt{\frac{2k+1}{2}},\qquad
 e_k(\lambda)=\sqrt{\frac{2k+1}{2}}\,L_k(\lambda),\quad k=0,\ldots,n.}
\]
\[
 e_0=\frac1{\sqrt2},\quad e_1=\sqrt{\frac32}\lambda,\quad
 e_2=\sqrt{\frac58}(3\lambda^2-1),\quad
 e_3=\sqrt{\frac78}(5\lambda^3-3\lambda).
\]
\newpage
\section{第 4 题}
证明教材定理 4.7 与 4.8。

\textbf{定理 4.7。} $n$ 阶矩阵 $A$ 为酉矩阵，当且仅当其 $n$ 个列（行）向量构成 $\C^n$ 的正交规范基。
\par\Needspace{4\baselineskip}\textbf{解答}\par

\[
 A=[a_1\ \cdots\ a_n],\qquad
 (A^*A)_{ij}=a_i^*a_j=\ip{a_i}{a_j}.
\]
\[
 A^*A=I\iff a_i^*a_j=\delta_{ij};\qquad
 A^*A=I\Longrightarrow A^{-1}=A^*\Longrightarrow AA^*=I.
\]
行向量情形由 $AA^*=I$ 同理得到；$n$ 个正交规范向量必为基。

\medskip
\textbf{定理 4.8。} 设 $\mathcal U\in L(V)$。下列条件等价：
\begin{enumerate}
\item $\mathcal U$ 是酉变换；
\item $\|\mathcal U\alpha\|=\|\alpha\|$, $\forall\alpha\in V$;
\item $\mathcal U$ 将一组正交规范基映为正交规范基；
\item $\mathcal U$ 在一组正交规范基下的矩阵为酉矩阵。
\end{enumerate}

\par\Needspace{4\baselineskip}\textbf{解答}\par

\[
 (1)\Longrightarrow(2):\quad
 \|\mathcal U\alpha\|^2=\ip{\mathcal U\alpha}{\mathcal U\alpha}=\ip\alpha\alpha.
\]
\begin{align*}
 (2)\Longrightarrow(1):\quad
 4\ip\alpha\beta
 &=\|\alpha+\beta\|^2-\|\alpha-\beta\|^2\\
 &\quad+i\|\alpha-i\beta\|^2-i\|\alpha+i\beta\|^2\\
 &=4\ip{\mathcal U\alpha}{\mathcal U\beta}.
\end{align*}
实数域只用上式的前两项。取正交规范基 $E=(e_1,\ldots,e_n)$，则
\[
 (1)\Longrightarrow(3):\quad
 \ip{\mathcal Ue_i}{\mathcal Ue_j}=\delta_{ij}.
\]
\[
 (3)\Longrightarrow(1):\quad
 \ip{\mathcal U\sum_i x_ie_i}{\mathcal U\sum_j y_je_j}
 =\sum_{i,j}\bar x_i y_j\delta_{ij}
 =\sum_i\bar x_i y_i=\ip{\sum_i x_ie_i}{\sum_j y_je_j}.
\]
\[
 C=[\mathcal U]_E:\quad
 \ip{\mathcal U\alpha}{\mathcal U\beta}=x^*C^*Cy,\qquad
 \ip\alpha\beta=x^*y
 \quad\Longrightarrow\quad (1)\iff C^*C=I\iff(4).
\]
\section{第 5 题}
设 $A\in\C^{n\times n}$，$r=\rank A$。证明 $A=A^*=A^2$ 当且仅当存在 $U\in\mathcal U^{n\times r}$ 使 $A=UU^*$，其中 $U^*U=I_r$。
\par\Needspace{4\baselineskip}\textbf{解答}\par

\textbf{$(\Leftarrow)$}
\[
 A^*=(UU^*)^*=UU^*=A,\qquad
 A^2=U(U^*U)U^*=UU^*=A,\qquad \rank A=\rank U=r.
\]
\textbf{$(\Rightarrow)$}
\[
 S=\im A,\quad T=\ker A,\qquad
 x=Ax+(I-A)x\in S+T,\qquad S\cap T=\{0\}.
\]
\[
 z\in T:\quad (Ay)^*z=y^*A^*z=y^*Az=0
 \quad\Longrightarrow\quad T=S^\perp.
\]
选 $S$ 的正交规范基，列成 $U=[u_1\ \cdots\ u_r]$。
\[
 U^*U=I_r,\qquad Ax\in S\Longrightarrow Ax=UU^*Ax,
 \qquad (I-A)x\in S^\perp\Longrightarrow U^*Ax=U^*x.
\]
\[
 Ax=UU^*x\quad(\forall x)\quad\Longrightarrow\quad\boxed{A=UU^*}.
\]
$r=0$ 时取 $n\times0$ 空矩阵 $U$，其乘积 $UU^*=0=A$。
\end{document}
